\section{Supplementary Evidence and Analysis Protocol}
\label{app:supplement}

\subsection{Traceable Examples in the Draft}

The main-text examples refer to the following archived run keys. Paths
below are relative to each run's workspace or trace root, not to an
author-specific machine. No firmware or exploit is executed by the file
inventory. The examples establish which retained artifacts motivated the
analysis; they do not replace complete outcome adjudication.

\paragraph{CVE-2023-26315 / claude-sonnet-4-6 / 1.}
The workspace includes \casefile{report.md},
\casefile{poc/final_exploit_run.txt}, extracted firmware contents, and
service logs including \casefile{logs/datacenter.log}. The trace directory
contains \casefile{metadata.json}, \casefile{session_messages.json}, and
container output. The report and demonstration identify a direct test of
the internal \texttt{plugincenter} component. The exact process/input/effect
chain and documented-entry scope require the checks in TODO CASE-01 and
RQ1-02. Values resembling credentials are not reproduced here.

\paragraph{CVE-2024-23624 / glm-5.2 / 2.}
The workspace includes \casefile{DAP-1650_v1.04B01.zip}, disassembly files,
and \casefile{ghidra_scripts/decompile.py}. The retained script contains
the misspelled \texttt{depiledFunction()} call. The corresponding trace
directory is empty. These observations do not establish the actual
exception sequence or stopping decision; see TODO CASE-02.

\paragraph{CVE-2020-27866 / claude-sonnet-4-6 / 1.}
The workspace includes \casefile{plan.md},
\casefile{vulnerable_server.py}, \casefile{exploit.py},
\casefile{server.log}, and \casefile{exploit_output.txt}. The plan and
server identify the simulation explicitly. The saved output reports nine
passing checks against the implemented server. The trace directory is
empty; consequently, the artifact set is not used to prove an exhaustive
search history or a transition after a timed rehosting failure.

\subsection{Initial Coding Fields}

The following fields define an initial analysis record. The codebook must
be refined and frozen through the pilot described in TODO COD-01.
An absent value is not silently mapped to a negative finding.

\begin{description}
\item[Run identity.] CVE, exact model route, trial, collection metadata,
and evidence availability.
\item[Activity.] Information gathering, acquisition, extraction,
component analysis, rehosting, or triggering; allow revisits and overlap.
\item[Action and feedback.] Observed command/tool call, relevant output,
timestamps or order, and references to the underlying evidence.
\item[Obstacle.] Unmet condition, direct evidence versus assertion,
whether repaired, and relation to termination if known.
\item[Response.] Repair, further search, repetition, target-preserving
alternative, simulation, return to original target, or stop; retain
multiple responses within a run.
\item[Outcome and scope.] Identity, execution and effect decisions;
documented entry versus internal component; unresolved prerequisites.
\item[Report.] Claimed result, disclosed limitations, simulation
disclaimer, and agreement with the observed execution scope.
\item[Adjudication.] Analyst decisions, independent overlap where
performed, disagreement, resolution, and remaining uncertainty.
\end{description}

\subsection{Sensitivity and Follow-Up Analyses}

The final reanalysis should report how findings change when restricted to
complete trajectories, when unresolved outcomes are retained explicitly,
and when related CVEs sharing an image or service are grouped. It should
also compare per-run results with three-trial group summaries. These
checks are pending rather than completed robustness results.
\todo{SUP-01}{Audit and complete the retained case/reference table,
recover the executed prompt, and publish coding rules, coverage strata,
and completed sensitivity results;
keep decisive methods and findings in the main body.}

Follow-up experiments are conditional on the final scientific claims.
Providing a ready original target addresses a different question from
asking for a clearer failure report. A second harness examines another
configuration. Any such experiment should retain matched cases, document
changed conditions, and report both positive and negative outcomes.
